% Figure 1 — one run is two phases, and they do not share a container.
%
% Left: the solve, which Inspect runs in a sandbox because the agent needs
% a shell. Right: the score, which runs on the host because the scorer
% starts the app's own stack. The only thing that crosses is the tree.
%
% Sources: AGENTS.md ("the bench owns the stack", "one sample, one stack"),
% src/bench/solves.py, src/bench/stack.py, specs/behavior.md §1-2.
\documentclass[tikz,border=3pt]{standalone}
\input{style}

\begin{document}
\begin{tikzpicture}[node distance=5mm and 6mm]

% ---------------------------------------------------------------- solve --
\node[file, text width=4.6cm] (wo)
  {\textbf{work order}\\the task type's contract\\ + the manifest's ask};
\node[device, text width=4.6cm, below=of wo] (agent)
  {\textbf{agent}\\one profile of \texttt{runs.yaml}};
\node[device, text width=4.6cm, below=of agent] (ws)
  {\textbf{workspace}\\the pinned tree at \texttt{/workspace},\\writable,
   caches warm};
\node[db, right=8mm of ws] (db) {db};
\node[note, below=2mm of db, text width=1.8cm, align=center] (dbnote)
  {MySQL 5.7,\\the app's own};
\node[file, text width=4.6cm, below=15mm of ws] (kept)
  {\textbf{out/solves/\textit{run}/}\\\texttt{record.json}\\
   \texttt{diff.patch} \quad \texttt{tree/}};

\draw[flow] (wo) -- node[lbl, right] {prompt} (agent);
\draw[flow] (agent) -- node[lbl, right] {edits} (ws);
\draw[link] (ws) -- (db);
\draw[flow] (ws) -- node[lbl, right] {cleanup hook} (kept);

\begin{scope}[on background layer]
  \node[sandbox, fit=(wo) (agent) (ws) (db) (dbnote) (kept), inner sep=4mm]
    (sandbox) {};
\end{scope}
\node[title, anchor=south, text width=7.6cm] at (sandbox.north)
  {1 \, solve — a sandbox, an agent, a tree};

% ---------------------------------------------------------------- score --
\node[device, text width=4.6cm, right=22mm of sandbox.east |- wo]
  (pw) {\textbf{Playwright}\\on the host};
\node[device, text width=4.6cm, below=of pw] (port)
  {the app's published port\\\texttt{APP\_HTTP\_PORT}};
\node[db, below=8mm of port] (db2) {db};
\node[note, below=3mm of db2, text width=3.6cm, align=center]
  {a fresh volume: empty\\at the start of every case};
\node[file, text width=4.6cm, below=19mm of db2] (delta)
  {\textbf{before} / \textbf{after} snapshots\\$\to$ delta against
   \texttt{expect.json}};

\draw[flow] (pw) -- (port);
\draw[flow] (port) -- (db2);
\draw[flow] (db2) -- (delta);

\begin{scope}[on background layer]
  \node[panel, fit=(pw) (port) (db2) (delta)] (scoring) {};
\end{scope}
\node[title, anchor=south, text width=7.4cm] at (scoring.north)
  {2 \, score — on the host, one stack per case};

% ---------------------------------------------------------- the handoff --
\draw[flow, line width=1pt, draw=teal]
  (kept.east) -- node[lbl, above, align=center]
  {\texttt{BENCH\_CHECKOUT}\\= the tree} (scoring.west |- kept);

\node[note, anchor=north west, text width=11.4cm]
  at ($(sandbox.south west) + (0,-7mm)$)
  {The driver carries exactly two things between the two phases: a run id,
   which names \texttt{out/solves/\textit{run}/}, and the tree that id
   names. A stale \texttt{BENCH\_CHECKOUT} is the one failure it cannot
   catch, which is why scoring derives it from the id rather than
   accepting it as input.};

\end{tikzpicture}
\end{document}
